% page 25 of the facsimile (image p-024 of the scan)
% block 165.1 x 242.0 mm ; left margin 36.9 mm
\pgc{15.240mm}{0.000mm}{
\l{\cel{54}17}
\l{}
\l{\sou{alciono}. (\sou{mitol.}) (\sou{en la epoki antiqua}). La alcedo, egardata}
\l{\cel{3}kom ucelo qua facas sua nesto sur maro kalma, e (da la}
\l{\cel{3}navigisti) kom ucelo bon-augura. - DEFIS.}
\l{}
\l{\sou{aleo}. I. Voyo quan facas bordumo de arbori, de arbusti folioza}
\l{\cel{3}o de herbi longa-stipa. - II. (\sou{arkeol.}) Paseyo tektoza,}
\l{\cel{3}unionuro ek du rangi inter-paralela de dolmeni. - DEFR.}
\l{}
\l{\sou{alegar}. (\sou{trans}.)Mencionar o citar kom autoritato, donar kom}
\l{\cel{3}justifikivo, kom punto di sua argumento. - DEFIS.}
\l{}
\l{\sou{alegorio}.\cel{2}Vorti figura qui prizentas a la spirito senco qua}
\l{\cel{3}esas celita sub la senco literala. - DEFIRS.}
\l{}
\l{\sou{alejar}. (\sou{trans.}) Igar min pezoza, Igar min kargoza. (\sou{anke metaf.})}
\l{\cel{3}- EFIS.}
\l{}
\l{\sou{aleno}. Puncilo quan uzas la shuifisti, la selifisti, por borar}
\l{\cel{3}trui en la ledro-peci quin li volas inter-sutar. - F.}
\l{}
\l{\sou{alerta}. Ago-rapida. - E.F.}
\l{}
\l{\sou{alexandrino}. verso de dek-e-du silabi, kun cezuro an-pos la}
\l{\cel{3}sisesma. - DEFIRS.}
\l{}
\l{\sou{alezar}. (\sou{trans}.) (\sou{teknol.})\cel{2}Polisar, glatigar la parti kava di}
\l{\cel{3}objekto borita o konkavigita. - FIS.}
\l{}
\l{\sou{alfabeto}. La serio de la literi qui reprezentas la soni ye qui}
\l{\cel{3}konsistas la vorti. - DEFIRS.}
\l{}
\l{\sou{alfenido}. Aloyuro de nikelo e de kupro. - DEFIRS.\cel{5}.}
\l{}
\l{\sou{algo}. (\sou{bot.}) Vejetanto kriptogama qua kreskas en aquo sen-sala}
\l{\cel{3}o saloza en la loki tre humida. - L.\cel{1}\sou{alga}.\cel{2}- DEFIRS.}
\l{}
\l{\sou{algebro}.\cel{2}(\sou{matem.}) Cienco di qua la skopo esas simpligar e solvar}
\l{\cel{3}la problemi pri la +grandori, reprezentante li per literi e}
\l{\cel{3}signi, e determinante la relati qui interligas la elementi}
\l{\cel{3}diversa di la problemo, per formulo quan onu povas aplikar ad}
\l{\cel{3}omna kazi simila ; e qua, preterirante la skopo, studias la}
\l{\cel{3}relati kom ipsajo, sen egardar lia origino, formi, kombinuri}
\l{\cel{3}e transformesi. - DEFIRS.}
\l{}
\l{\sou{algoritmo}.\cel{2}(\sou{matem.}) Tipo di noto-sistemo qua esas propra di}
\l{\cel{3}genero partikulara di kalkulado. - DEFIRS.}
\l{}
\l{\sou{alio}. (\sou{bot.})\cel{2}Varietato di planti liliacea, di qua la bulbo}
\l{\cel{3}odoras e saporas forte e pikante : shaloto, oniono, cibolo,}
\l{\cel{3}ciboleto. - L.\cel{1}\sou{allium sativum}. - FISL.}
\l{}
\l{\sou{aliancar}. (\sou{trans., kun)}.\cel{2}Unionar personi, familii, per mariajo.}
\l{\cel{3}- DEFIS. (Ta vorto ne esas sinonimo di\cel{1}\sou{mariajar}\cel{1}:\cel{1}\sou{mariajo}}
\l{\cel{3}relatas nur la gespozi ;\cel{1}\sou{aliancar}\cel{1}relatas lia familii.)}
}
